% ============================================================================
% Fichier  : partie4/P04_C13_gouvernance_projets_prestataires.tex
% Titre    : Gouvernance, Projets Numériques et Prestataires
% Auteur   : MINKA MI NGUIDJOÏ Thierry Emmanuel
% Version  : 0.1 -- novembre 2026
% ============================================================================
% Sources : INF3226 C17, C19, C24 (resserrés) ; COBIT 2019 (gouvernance,
% projets) ; GUID 5100 ; loi n° 2024/017 (décrite sans citation d'article).
% Faits e-RECETTES établis ici : marché de 640 000 000 FCFA attribué à
% NUMÉRIKA SARL en mars 2024 ; mise en service le 2 janvier 2025 ; PV de
% réception du 15 janvier 2025 signé avec 47 anomalies ouvertes, dont 9
% bloquantes ; aucun schéma directeur ; pas de clause de réversibilité ni
% de propriété des données ; hébergement hors du pays ; compte du
% prestataire avec droits d'administration en production ; 12 400
% contribuables enregistrés (données personnelles).
% ============================================================================

\chapter{Gouvernance, Projets Numériques et Prestataires}
\label{chap:gouvernance}
\chaptermark{Gouvernance, projets, prestataires}

\epigraph{%
    « On peut confier un service~; on ne confie pas sa responsabilité. »%
}{Principe du cours}

\niveauChapitre{Intermédiaire}{Chapitre~\ref{chap:sci}.}

\begin{boxobjectifs}
\begin{itemize}
    \item Évaluer la gouvernance de la fonction informatique d'une
          administration~: schéma directeur, organisation, pilotage.
    \item Auditer un projet numérique et le marché qui le porte, de
          l'expression du besoin à la réception.
    \item Évaluer les contrats avec les prestataires~: propriété des
          données, réversibilité, accès, niveaux de service.
    \item Repérer les risques de l'hébergement externe.
    \item Situer les obligations de protection des données personnelles.
\end{itemize}
\end{boxobjectifs}

\bigskip

% ============================================================================
\section*{Pourquoi ce chapitre}
% ============================================================================

Les déficiences des chapitres précédents ont souvent une cause commune,
plus haut~: un système acheté sans besoin écrit, reçu sans tests, confié
à un prestataire sans contrat solide, piloté par personne. Ce chapitre
remonte à cette cause, sur le dossier e-RECETTES.

\separateur

% ============================================================================
\section{La gouvernance de la fonction informatique}
\label{sec:c13-gouvernance}
% ============================================================================

\begin{tableaucours}{L{3.2cm}L{4.4cm}L{4.6cm}}{%
    \tableauHeader{Élément} &
    \tableauHeader{Question} &
    \tableauHeader{Pièce}
}
Schéma directeur & Les projets répondent-ils à des objectifs écrits~? &
Schéma directeur adopté, suivi annuel \\
Organisation & Qui décide, qui exécute, qui contrôle~? & Organigramme,
décisions, propriétaires d'application \\
Pilotage & Un comité suit-il les projets et les risques~? & Comptes rendus
\\
Budget & Le coût complet du SI est-il connu~? & Budget informatique,
contrats \\
Compétences & L'entité peut-elle contrôler ses prestataires~? & Effectifs,
formations \\
\end{tableaucours}
\vspace{-0.8em}
\begin{center}{\small\itshape Tableau~13.1. Les éléments de la gouvernance
de la fonction informatique}\end{center}

COBIT~2019 sépare la \emph{gouvernance} (évaluer, diriger, surveiller),
qui relève de la direction, de la \emph{gestion}, qui relève des
services~\cite{isaca2018cobitintro}. Dans une administration, la
gouvernance revient à l'ordonnateur et à la direction~; elle ne peut pas
être laissée au prestataire.

\separateur

% ============================================================================
\section{Auditer un projet numérique}
\label{sec:c13-projet}
% ============================================================================

\subsection{Le cycle du projet et ses contrôles}

\begin{tableaucours}{L{2.6cm}L{5.0cm}L{4.6cm}}{%
    \tableauHeader{Phase} &
    \tableauHeader{Contrôle attendu} &
    \tableauHeader{Pièce}
}
Besoin & Besoin écrit par le métier, rattaché au schéma directeur & Note
de besoin, étude d'opportunité \\
Marché & Cahier des charges qui inclut les contrôles internes attendus,
la sécurité, la réversibilité & Dossier d'appel d'offres, contrat \\
Réalisation & Suivi par un comité, avec les risques & Comptes rendus \\
Recette & Tests par le métier, anomalies bloquantes levées avant
réception & Cahier de recette, PV de réception \\
Mise en service & Reprise des données contrôlée, formation & PV de reprise,
totaux de contrôle \\
\end{tableaucours}
\vspace{-0.8em}
\begin{center}{\small\itshape Tableau~13.2. Les contrôles du cycle d'un
projet}\end{center}

\begin{boxscenario}{Dossier e-RECETTES. Une réception trop tôt}
\textbf{Constatation.} Le marché de 640\,000\,000~FCFA a été attribué à
NUMÉRIKA~SARL en mars~2024. e-RECETTES est en service depuis le
2~janvier~2025. Le procès-verbal de réception, signé le 15~janvier~2025,
mentionne 47 anomalies ouvertes, dont 9 bloquantes. Le cahier des charges
ne prévoit ni contrôle des annulations, ni numérotation continue des
quittances. Il n'existe pas de schéma directeur.

\medskip
\textbf{Analyse} (confiance élevée). Les déficiences constatées aux
Chapitres~\ref{chap:exploitation} et~\ref{chap:applications} étaient
inscrites dans le projet~: ce qui n'est pas demandé au cahier des
charges n'est pas livré, et ce qui est reçu avec des anomalies
bloquantes est payé sans être conforme.
\end{boxscenario}

\begin{boxalert}
\textbf{Violation typique.} Prononcer la réception d'un système pour
respecter le calendrier budgétaire, avec des anomalies bloquantes
ouvertes.

\smallskip
\textbf{Règle~:} une réception avec réserves exige un plan de levée daté
et une retenue sur paiement~; une réception sans réserve d'un système
défaillant engage la responsabilité de celui qui la signe.
\end{boxalert}

\separateur

% ============================================================================
\section{Les prestataires et l'hébergement externe}
\label{sec:c13-prestataires}
% ============================================================================

\begin{tableaucours}{L{3.4cm}L{8.8cm}}{%
    \tableauHeader{Clause} &
    \tableauHeader{Ce que l'auditeur vérifie}
}
Propriété des données & Les données appartiennent à l'administration, qui
peut les obtenir à tout moment dans un format exploitable \\
Réversibilité & Les conditions de reprise du service par l'administration
ou un autre prestataire sont écrites \\
Accès du prestataire & Ses accès en production sont nominatifs, limités,
journalisés, et ouverts sur demande \\
Niveaux de service & Disponibilité, délais de correction, pénalités \\
Localisation & Lieu d'hébergement connu et compatible avec les obligations
légales \\
Droit d'audit & L'administration et ses contrôleurs peuvent auditer le
prestataire \\
\end{tableaucours}
\vspace{-0.8em}
\begin{center}{\small\itshape Tableau~13.3. Les clauses à vérifier dans un
contrat de prestation informatique}\end{center}

\begin{boxscenario}{Dossier e-RECETTES. Le contrat de NUMÉRIKA}
\textbf{Constatation.} Le contrat ne contient ni clause de propriété des
données, ni clause de réversibilité, ni droit d'audit. Les données sont
hébergées hors du pays. Le prestataire dispose d'un compte
d'administration permanent en production (Chapitre~\ref{chap:exploitation}).

\medskip
\textbf{Analyse} (confiance élevée). La recette des finances dépend
entièrement de NUMÉRIKA pour accéder à ses propres données. Le décret
n°~2013/287 oblige les responsables de l'entité auditée à présenter les
pièces~; s'ils ne peuvent pas les obtenir de leur prestataire, c'est
l'entité qui est en défaut.
\end{boxscenario}

\separateur

% ============================================================================
\section{Les données personnelles}
\label{sec:c13-donnees}
% ============================================================================

Les systèmes financiers publics traitent des données personnelles~: agents
(paie), contribuables, usagers, fournisseurs personnes physiques. La loi
n°~2024/017 du 23~décembre~2024 encadre ces traitements au Cameroun. Le
cours n'en cite pas les articles~; l'auditeur les lit avant de s'y
référer. Il vérifie au moins que les traitements sont recensés, que
l'accès aux données est limité aux personnes qui en ont besoin, et que
les transferts vers un hébergeur, en particulier hors du pays, sont
encadrés.

\begin{boxinfo}
\textbf{Les extractions de l'auditeur sont aussi des données
personnelles.} Le fichier de paie de la CAPNO, avec ses 4\,312 noms et
comptes bancaires, doit être protégé par le vérificateur avec le même
soin qu'il l'exige de l'entité (Chapitre~\ref{chap:normes}).
\end{boxinfo}

\begin{boxlecture}{Gouvernance et projet e-RECETTES}
\textbf{Constatation.} Sur huit contrôles attendus (schéma directeur,
cahier des charges complet, comité de suivi, réception sans anomalie
bloquante, propriété des données, réversibilité, accès du prestataire
encadré, droit d'audit), un seul fonctionne~: un comité de suivi s'est
réuni trois fois pendant le projet.

\textbf{Analyse} (confiance élevée). La cause première des déficiences
d'e-RECETTES est dans la gouvernance du projet.

\medskip
\hfill\conformite{1}{8}{0}
\end{boxlecture}

\begin{boxmetier}{Vérificateur}
Pour tout système acquis par marché, demandez le cahier des charges et le
procès-verbal de réception. Comparez les contrôles demandés aux contrôles
livrés~: l'écart explique souvent les déficiences.
\end{boxmetier}

\begin{boxmetier}{Inspecteur}
Une réception signée avec des anomalies bloquantes ouvertes est un fait
précis, daté et signé. Elle mérite une question écrite à son signataire.
\end{boxmetier}

\begin{boxmetier}{Contrôleur financier}
Avant de viser le paiement d'un solde de marché informatique, demandez le
procès-verbal de réception et la liste des réserves levées.
\end{boxmetier}

\separateur

\begin{boxresume}
\begin{itemize}
    \item \textbf{La gouvernance reste à l'administration}~: schéma
          directeur, propriétaires, comité, budget, compétences.
    \item \textbf{Le projet porte les contrôles}~: ce qui n'est pas au
          cahier des charges n'est pas livré.
    \item \textbf{La réception engage}~; des anomalies bloquantes ouvertes
          imposent réserves, plan et retenue.
    \item \textbf{Six clauses de contrat}~: données, réversibilité, accès,
          niveaux de service, localisation, droit d'audit.
    \item \textbf{Les données personnelles} des systèmes financiers sont
          encadrées par la loi n°~2024/017, y compris chez l'auditeur.
\end{itemize}
\end{boxresume}

\begin{boxexo}[title={\ding{45}\quad Exercice~13.1. Gouvernance ou gestion~? \niveaubase}]
Classez~: (a) adopter un schéma directeur~; (b) faire les sauvegardes~;
(c) décider d'externaliser l'hébergement~; (d) attribuer un profil.
\end{boxexo}

\begin{boxexo}[title={\ding{45}\quad Exercice~13.2. Un cahier des charges \niveaumoyen}]
Rédigez les six exigences de contrôle interne que le cahier des charges
d'e-RECETTES aurait dû contenir, à partir des
Chapitres~\ref{chap:acces} à~\ref{chap:applications}.
\end{boxexo}

\begin{boxexo}[title={\ding{45}\quad Exercice~13.3. Le contrat \niveaumoyen}]
Rédigez la recommandation sur le contrat de NUMÉRIKA, avec les clauses à
ajouter par avenant et l'ordre de priorité.
\end{boxexo}

\begin{boxexo}[title={\ding{45}\quad Exercice~13.4. Les causes \niveauexpert}]
Reliez chaque déficience d'e-RECETTES relevée aux Chapitres~\ref{chap:exploitation}
et~\ref{chap:applications} à une défaillance du projet ou du contrat, et
dites quelle recommandation de gouvernance les traiterait toutes.
\end{boxexo}

\begin{boxinfo}
\textbf{Transition.} Toutes les pièces sont en place. Le
Chapitre~\ref{chap:cas-integrateur} les réunit dans une mission
complète, conduite en équipe en deux heures.
\end{boxinfo}
